%%TCC -- Roteiro de apresentação (defesa)
%%Predição de Evasão Estudantil no Curso de SI da UFVJM
%%Estrutura sugerida pelo orientador: Objetivo -> Hipóteses -> Metodologia -> Resultados -> Dashboard (ao final)
%%Compilar com: pdflatex apresentacao.tex  (ou latexmk -pdf apresentacao.tex)

\documentclass[aspectratio=169,11pt,brazil]{beamer}

\usepackage[utf8]{inputenc}
\usepackage[T1]{fontenc}
\usepackage[brazil]{babel}
\usepackage{mathptmx}          % Times New Roman, consistente com o TCC
\usepackage{graphicx}
\usepackage{booktabs}
\usepackage{xcolor}

% --- Identidade visual ---------------------------------------------------
\definecolor{ufvjmgreen}{RGB}{0,96,64}
\definecolor{ufvjmgray}{RGB}{90,90,90}

\usetheme{Madrid}
\usecolortheme[named=ufvjmgreen]{structure}
\setbeamercolor{palette primary}{bg=ufvjmgreen,fg=white}
\setbeamercolor{palette secondary}{bg=ufvjmgreen!80,fg=white}
\setbeamercolor{palette tertiary}{bg=ufvjmgreen!60,fg=white}
\setbeamercolor{title}{bg=ufvjmgreen,fg=white}
\setbeamercolor{subtitle}{fg=white}
\setbeamercolor{frametitle}{fg=ufvjmgreen}
\setbeamertemplate{navigation symbols}{}
\setbeamertemplate{caption}{\insertcaption}
\setbeamerfont{frametitle}{series=\bfseries}
\graphicspath{{reports/figures/}{./}}

% --- Dados do trabalho ---------------------------------------------------
\title[Predição de Evasão no SI/UFVJM]{Predição de evasão estudantil no curso de\\Sistemas de Informação da UFVJM}
\subtitle{Uma Análise por Ciência de Dados e Aprendizado de Máquina}
\author[Josiney M. B. Junior]{Josiney Mafra Barbosa Junior}
\institute[UFVJM]{Orientador: Prof. Dr. Alessandro Vivas\\Universidade Federal dos Vales do Jequitinhonha e Mucuri}
\date{2026}

\begin{document}

% ======================================================================
% BLOCO 1 -- ABERTURA E CONTEXTO  (~2 min)
% ======================================================================
\begin{frame}[plain]
  \titlepage
\end{frame}

\begin{frame}{Contexto: por que estudar a evasão?}
  \begin{columns}[T]
    \begin{column}{0.55\textwidth}
      \begin{itemize}
        \item Evasão na grande área de \textbf{Computação} no Brasil ultrapassa \textbf{60\%}.
        \item Fenômeno \textbf{multifatorial}: condições acadêmicas, socioeconômicas, institucionais e pessoais.
        \item Custos em três níveis: \textbf{estudante}, \textbf{universidade} e \textbf{sociedade}.
      \end{itemize}
    \end{column}
    \begin{column}{0.45\textwidth}
      \begin{block}{Oportunidade analítica}
        Cada estudante deixa um \textbf{rastro digital} nos sistemas institucionais (matrículas, notas, frequência).
        \medskip

        A Ciência de Dados permite \textbf{antecipar} o desfecho e apoiar um \textbf{Sistema de Alerta Precoce (SAP)}.
      \end{block}
    \end{column}
  \end{columns}
\end{frame}

% ======================================================================
% BLOCO 2 -- OBJETIVO  (~1,5 min)
% ======================================================================
\section{Objetivo}

\begin{frame}{Objetivo}
  \begin{block}{Objetivo geral}
    Aplicar técnicas de \textbf{Ciência de Dados e Aprendizado de Máquina} para \textbf{identificar e analisar os fatores preditivos da evasão e da retenção prolongada} dos estudantes do curso de Sistemas de Informação da UFVJM.
  \end{block}
  \begin{block}{Objetivos específicos}
    \begin{enumerate}
      \item \textbf{Explorar} e preparar os dados, garantindo qualidade e consistência.
      \item \textbf{Modelar} a evasão (Regressão Logística, \emph{Random Forest}, \emph{XGBoost}) e a dimensão temporal (Kaplan-Meier e Cox).
      \item \textbf{Avaliar} o desempenho com métricas sensíveis ao desbalanceamento e identificar os atributos de maior impacto.
    \end{enumerate}
  \end{block}
  \centering\small Produto final: subsídio para um \textbf{Sistema de Alerta Precoce}.
\end{frame}

% ======================================================================
% BLOCO 3 -- HIPÓTESES  (~2 min)
% ======================================================================
\section{Hipóteses}

\begin{frame}{Hipóteses de pesquisa}
  \begin{itemize}
    \item[\textbf{H1}] A \textbf{reprovação por frequência} é um dos preditores mais relevantes da evasão.
    \item[\textbf{H2}] O \textbf{ciclo de Matemática} (MAT001--MAT003) é o filtro estrutural dominante, acima das disciplinas iniciais de Computação.
    \item[\textbf{H3}] O \textbf{CRA} dos primeiros semestres é indicador precoce de risco.
    \item[\textbf{H4}] O \textbf{tempo prolongado sem progressão} curricular eleva o risco de evasão.
    \item[\textbf{H5}] \textbf{\emph{Ensembles} de árvore} (RF, XGBoost) superam a Regressão Logística em métricas sensíveis ao desbalanceamento.
  \end{itemize}
  \vfill
  \centering\small Cada hipótese é retomada nos \textbf{Resultados} com veredito de confirmação.
\end{frame}

% ======================================================================
% BLOCO 4 -- METODOLOGIA  (~4 min)
% ======================================================================
\section{Metodologia}

\begin{frame}{Metodologia: CRISP-DM como fio condutor}
  \begin{columns}[c]
    \begin{column}{0.52\textwidth}
      \begin{itemize}
        \item Abordagem \textbf{quantitativa, orientada a dados}.
        \item Organizada nas \textbf{6 fases} do CRISP-DM.
        \item Cada fase materializada em artefatos: \texttt{src/}, \emph{notebooks} e painel \emph{Streamlit}.
      \end{itemize}
    \end{column}
    \begin{column}{0.48\textwidth}
      \centering
      \includegraphics[width=\linewidth,height=0.72\textheight,keepaspectratio]{fig_crispdm_ciclo.pdf}
    \end{column}
  \end{columns}
\end{frame}

\begin{frame}{Dados e preparação}
  \begin{columns}[T]
    \begin{column}{0.5\textwidth}
      \textbf{Base institucional (SIE/UFVJM)} --- 3 tabelas integradas:
      \begin{itemize}
        \item \textbf{1.119 estudantes}
        \item \textbf{37.925 registros disciplinares}
      \end{itemize}
      \medskip
      \textbf{Decisões-chave}
      \begin{itemize}
        \item Recorte: ingressantes até 2024/1.
        \item Exclusão de \texttt{COM038} (TCC).
        \item Tabela analítica: \textbf{uma linha por aluno}.
      \end{itemize}
    \end{column}
    \begin{column}{0.5\textwidth}
      \begin{block}{\emph{Features} ancoradas nas hipóteses}
        \small
        CRA e CRA médio, taxa de aprovação, reprovações por nota/frequência, tentativas em MAT001--003, trancamentos, variância das notas.
      \end{block}
      \begin{alertblock}{Ressalva metodológica}
        \small As \emph{features} são \textbf{cumulativas}.
      \end{alertblock}
    \end{column}
  \end{columns}
\end{frame}

\begin{frame}{Modelagem e avaliação}
  \begin{columns}[T]
    \begin{column}{0.5\textwidth}
      \begin{block}{Duas frentes complementares}
        \begin{itemize}
          \item \textbf{Sobrevivência} (Kaplan-Meier + Cox): \emph{quando} a evasão ocorre.
          \item \textbf{Classificação} (LogReg, RF, XGBoost): \emph{quem} tende a evadir.
        \end{itemize}
      \end{block}
    \end{column}
    \begin{column}{0.5\textwidth}
      \begin{block}{Avaliação}
        \begin{itemize}
          \item \textbf{F1, \emph{recall}, precisão, AUC} (acurácia engana em dados desbalanceados).
          \item Validação cruzada estratificada (5 \emph{folds}) + \emph{hold-out}.
          \item Interpretabilidade via \textbf{SHAP}.
          \item Reprodutibilidade: \texttt{random\_state = 42}.
        \end{itemize}
      \end{block}
    \end{column}
  \end{columns}
\end{frame}

% ======================================================================
% BLOCO 5 -- RESULTADOS  (~7 min)
% ======================================================================
\section{Resultados}

\begin{frame}{Caracterização da coorte}
  \begin{columns}[c]
    \begin{column}{0.46\textwidth}
      \begin{itemize}
        \item 587 \textbf{Evadidos} (52,5\%), 260 \textbf{Concluintes} (23,2\%), 271 \textbf{Ativos} (24,2\%).
        \item Taxa de evasão (desfecho consolidado): \[ \frac{587}{587+260} = \textbf{69,3\%} \]
        \item Compatível com a literatura ($>$60\% em Computação).
      \end{itemize}
      \footnotesize\color{ufvjmgray} Cautela: coortes recentes sofrem censura à direita.
    \end{column}
    \begin{column}{0.54\textwidth}
      \centering
      \includegraphics[width=\textwidth]{fig_situacao_e_evasao_por_ano.pdf}
    \end{column}
  \end{columns}
\end{frame}

\begin{frame}{H3 --- O CRA separa os grupos}
  \begin{columns}[c]
    \begin{column}{0.48\textwidth}
      \begin{itemize}
        \item Concluintes: mediana de CRA $\approx$ \textbf{73}.
        \item Evadidos: mediana $\approx$ \textbf{45}.
        \item Ativos: comportamento próximo ao dos concluintes.
      \end{itemize}
      \medskip
      \begin{block}{Leitura}
        Primeira evidência descritiva forte a favor de \textbf{H3}.
      \end{block}
    \end{column}
    \begin{column}{0.52\textwidth}
      \centering
      \includegraphics[width=\textwidth]{fig_cra_por_situacao.pdf}
    \end{column}
  \end{columns}
\end{frame}

\begin{frame}{H2 --- Disciplinas-gargalo}
  \begin{columns}[c]
    \begin{column}{0.5\textwidth}
      \centering
      \includegraphics[width=\textwidth]{fig_top_gargalos.pdf}
    \end{column}
    \begin{column}{0.5\textwidth}
      \begin{itemize}
        \item Ciclo de Matemática nas 3 primeiras posições: \textbf{MAT003 63\%}, \textbf{MAT002 57\%}, \textbf{MAT001 53\%}.
        \item \texttt{COM001} em 5º (48\%) e \texttt{COM006} em 6º (45\%).
        \item \texttt{MAT001}: maior \textbf{evasão subsequente} (52\%).
      \end{itemize}
      \begin{block}{Veredito}
        \textbf{H2 confirmada}: Matemática supera Computação.
      \end{block}
    \end{column}
  \end{columns}
\end{frame}

\begin{frame}{H4 --- Análise de sobrevivência}
  \begin{columns}[c]
    \begin{column}{0.5\textwidth}
      \centering
      \includegraphics[width=\textwidth]{fig_sobrevivencia_geral.pdf}
    \end{column}
    \begin{column}{0.5\textwidth}
      \begin{itemize}
        \item \textbf{11\%} evadem já no 1º semestre; $\sim$46\% até o 9º (tempo regular).
        \item Mediana de permanência $\approx$ \textbf{11 semestres}.
        \item Estratificação por CRA: \textbf{log-rank $p<0{,}001$}.
        \item Modelo de Cox: índice de concordância \textbf{0,913}.
      \end{itemize}
      \footnotesize\color{ufvjmgray} HR $<$ 1 das reprovações cumulativas (viés de seleção).
    \end{column}
  \end{columns}
\end{frame}

\begin{frame}{H5 --- Desempenho dos modelos}
  \centering
  \begin{table}
    \small
    \begin{tabular}{lcccc}
      \toprule
      \textbf{Modelo} & \textbf{F1} & \textbf{AUC} & \textbf{\emph{Recall}} & \textbf{Precisão} \\
      \midrule
      Regressão Logística & 0,943 & 0,976 & 0,908 & \textbf{0,980} \\
      \emph{Random Forest} & 0,953 & 0,981 & 0,944 & 0,962 \\
      \emph{XGBoost} & \textbf{0,957} & \textbf{0,982} & \textbf{0,947} & 0,967 \\
      \bottomrule
    \end{tabular}
    \caption*{\footnotesize Médias na validação cruzada estratificada (5 \emph{folds}).}
  \end{table}
  \begin{columns}[T]
    \begin{column}{0.5\textwidth}
      \begin{itemize}
        \item \emph{Hold-out}: RF melhor F1 (\textbf{0,966}); \emph{recall} RF/XGB $=0{,}949$.
        \item \emph{Recall} importa num SAP: falso negativo custa mais.
      \end{itemize}
    \end{column}
    \begin{column}{0.5\textwidth}
      \begin{block}{Veredito}
        \textbf{H5 confirmada}: \emph{ensembles} superam o \emph{baseline} em F1 e \emph{recall}.
      \end{block}
    \end{column}
  \end{columns}
\end{frame}

\begin{frame}{Interpretabilidade (SHAP) e achados emergentes}
  \begin{columns}[c]
    \begin{column}{0.5\textwidth}
      \centering
      \includegraphics[width=\textwidth]{fig_importancia_rf.pdf}
    \end{column}
    \begin{column}{0.5\textwidth}
      \begin{itemize}
        \item Topo: \texttt{taxa\_aprovacao}, \texttt{cra\_medio}, \texttt{reprovacoes\_nota}, \texttt{cra}.
        \item \texttt{tentativas\_mat002/003} em 4º--6º $\Rightarrow$ \textbf{H2} em nível individual.
        \item \texttt{reprovacoes\_frequencia} em 8º $\Rightarrow$ \textbf{H1 parcial}.
        \item \texttt{tipoingresso\_Vestibular} em 12º $\Rightarrow$ discussão de \textbf{viés algorítmico}.
      \end{itemize}
    \end{column}
  \end{columns}
\end{frame}

\begin{frame}{Placar das hipóteses}
  \centering
  \begin{table}
    \renewcommand{\arraystretch}{1.3}
    \begin{tabular}{clc}
      \toprule
      & \textbf{Hipótese} & \textbf{Veredito} \\
      \midrule
      H1 & Reprovação por frequência & \color{orange}\textbf{Parcial} \\
      H2 & Ciclo de Matemática como filtro dominante & \color{ufvjmgreen}\textbf{Confirmada} \\
      H3 & CRA como indicador precoce & \color{ufvjmgreen}\textbf{Confirmada} \\
      H4 & Tempo sem progressão & \color{orange}\textbf{Parcial} \\
      H5 & \emph{Ensembles} de árvore superam \emph{baseline} & \color{ufvjmgreen}\textbf{Confirmada} \\
      \bottomrule
    \end{tabular}
  \end{table}
\end{frame}

% ======================================================================
% BLOCO 6 -- DASHBOARD (AO FINAL, conforme orientador)  (~2,5 min)
% ======================================================================
\section{Dashboard}

\begin{frame}{Dashboard --- demonstração ao vivo}
  \begin{columns}[T]
    \begin{column}{0.5\textwidth}
      \textbf{Painel interativo em \emph{Streamlit}}
      \begin{itemize}
        \item Visão temporal da evasão por coorte.
        \item Mapeamento das disciplinas-gargalo.
        \item \textbf{Semáforo de risco individual} com explicação.
      \end{itemize}
      \medskip
      Cada alerta vem com \textbf{justificativa quantificada}. A decisão final é da \textbf{coordenação}.
    \end{column}
    \begin{column}{0.5\textwidth}
      \begin{block}{Enfrenta o ``problema da última milha''}
        Traduz o rigor técnico em \textbf{instrumento acionável} pela gestão.
      \end{block}
      % \includegraphics[width=\textwidth]{dashboard_screenshot.png} % inserir print do painel
    \end{column}
  \end{columns}
\end{frame}

% ======================================================================
% BLOCO 7 -- FECHAMENTO  (~1 min)
% ======================================================================
\section{Considerações finais}

\begin{frame}{Contribuições, limitações e trabalhos futuros}
  \begin{columns}[T]
    \begin{column}{0.34\textwidth}
      \begin{block}{Contribuições}
        \footnotesize
        \begin{itemize}
          \item Painel acionável (SAP).
        \end{itemize}
      \end{block}
    \end{column}
    \begin{column}{0.34\textwidth}
      \begin{block}{Limitação central}
        \footnotesize
        Modelo é \textbf{diagnóstico retrospectivo}, não predição prospectiva pura (\emph{features} cumulativas).
      \end{block}
    \end{column}
  \end{columns}
\end{frame}

\begin{frame}[plain]
  \centering
  \vfill
  {\Huge\color{ufvjmgreen}\textbf{Obrigado!}}
  \vfill
  {\large Josiney Mafra Barbosa Junior}\\[0.3em]
  {\normalsize Orientador: Prof. Dr. Alessandro Vivas}\\[0.3em]
  {\small Universidade Federal dos Vales do Jequitinhonha e Mucuri}
  \vfill
  {\normalsize\itshape À disposição para perguntas.}
  \vfill
\end{frame}

\end{document}
